\chapter{Testing, Evaluation, and Future Work}
\label{chap:testing}
\label{chap:conclusion}

This chapter evaluates the current implementation. The evaluation focuses on
correctness of the implemented compiler pipeline, example coverage, snapshot
testing, and the limitations that remain. It also presents future work that
would move \epl{} closer to a production language.

\section{Testing Strategy}

The project uses a lightweight test harness in the \texttt{epl\_test} crate. The
harness recursively scans provided directories for \texttt{.epl} files. For each
source file, it performs two checks.

First, it runs the compiler in \texttt{ir\_tree} mode and compares the output to
a neighboring \texttt{.epl.ir\_tree} snapshot. If the snapshot is missing, the
harness writes a new one and marks the test as new. If the snapshot exists but
differs, the test fails. This catches changes to semantic lowering, type
checking, compile-time evaluation, and tree cleanup.

Second, it runs the compiler in \texttt{ir} mode as a smoke test. The current
harness does not snapshot lower IR output, but it verifies that lower IR
generation succeeds for every example and regression program.

\section{Verification Result}

The implementation was built and tested in the project workspace. The compiler
build command completed successfully:

\begin{lstlisting}[caption={Compiler build command}]
cargo build
\end{lstlisting}

The build produced one platform linker warning about a macOS library deployment
target, but the compiler binary was produced successfully. The test harness was
then run with:

\begin{lstlisting}[caption={Test harness command}]
cargo run -p epl_test ./target/debug/epl examples tests
\end{lstlisting}

The result was:

\begin{lstlisting}[caption={Test harness summary}]
DONE (24 passed, 0 failed, 0 new)
\end{lstlisting}

The 24 checks come from 12 \texttt{.epl} source files. Each file is tested once
against its \irtree{} snapshot and once as a lower IR smoke test.

\section{Example Coverage}

The example and test programs cover the main implemented language features.

{\footnotesize
\begin{longtable}{p{0.44\textwidth}p{0.38\textwidth}}
  \caption{Example and regression source files} \label{tab:examples} \\
  \toprule
  Source file & Feature coverage \\
  \midrule
  \endfirsthead
  \toprule
  Source file & Feature coverage \\
  \midrule
  \endhead
  \texttt{examples/hello\_world.epl} & External variadic declaration and string literal call to \texttt{printf}. \\
  \texttt{examples/fibonacci.epl} & Recursive function, iterative function, \texttt{while}, \texttt{for}, arithmetic, and external printing. \\
  \texttt{examples/fibonacci\_array.epl} & Fixed-size arrays, array mutation, array return, and array printing. \\
  \texttt{examples/comptime\_primes.epl} & \texttt{comptime}, pure function calls, nested loops, arrays, casts, and compile-time mutation. \\
  \texttt{examples/struct.epl} & Struct definition, named and unnamed initializers, by-value struct calls, field access, and compile-time struct computation. \\
  \texttt{examples/pointers.epl} & Address-of, typed pointer dereference, and pointer mutation. \\
  \texttt{examples/elseif.epl} & Nested \texttt{else if} expressions and branch result typing. \\
  \texttt{examples/exit.epl} & External never-returning function declaration. \\
  \texttt{tests/dead\_elimination.epl} & Tree cleanup after \texttt{return}. \\
  \texttt{tests/useless\_nested\_blocks.epl} & Block simplification. \\
  \texttt{tests/redundant\_deref\_ops.epl} & Simplification of redundant address/dereference pairs. \\
  \texttt{tests/i32\_min.epl} & Integer literal edge case. \\
  \bottomrule
\end{longtable}
}

This coverage is useful, but it is not exhaustive. It tests representative
successful programs and a few optimization cases. It does not yet include many
negative tests for diagnostics, ABI edge cases, malformed programs, or
platform-specific behavior.

\section{Snapshot Testing}

Snapshot testing is particularly effective for this compiler because \irtree{}
is readable. When a lowering rule changes, the generated snapshot changes in a
way that can be inspected by the developer. For example, the
\texttt{redundant\_deref\_ops.epl} test verifies that \texttt{(\&a).*} becomes
a load from \texttt{a} and that \texttt{\&(a\_ptr.*)} becomes a load of
\texttt{a\_ptr}. The snapshot exposes the optimizer's behavior directly.

Snapshot testing also has risks. A developer can accidentally accept an
incorrect new snapshot. The snapshot format is not a formal language
specification. It is a debugging representation. For this project, however, the
benefits are significant because the snapshots document the compiler pipeline
and catch regressions cheaply.

\section{Manual Inspection}

The command-line modes support manual inspection beyond the automated harness.
For example:

\begin{lstlisting}[caption={Inspecting generated LLVM IR}]
cargo run -- llvm-ir examples/hello_world.epl
\end{lstlisting}

The \texttt{cfg} command can be combined with Graphviz:

\begin{lstlisting}[caption={Generating a CFG graph}]
cargo run -- cfg examples/fibonacci.epl > /tmp/cfg.dot
dot -Tsvg /tmp/cfg.dot -o /tmp/cfg.svg
\end{lstlisting}

These inspection modes are valuable during compiler development because they
make each stage visible. They are also useful for explaining the project during
presentation or assessment.

\section{Correctness Discussion}

The current tests show that the implemented examples lower successfully and
that \irtree{} output remains stable. They do not prove full compiler
correctness. Compiler correctness would require a much broader strategy:
negative tests, randomized programs, differential testing against a reference
interpreter, run-time output checks, LLVM IR checks, and tests on multiple
platforms.

Even so, the current verification is meaningful. It exercises the full front end
and semantic pipeline for every example. It exercises compile-time evaluation
through the prime and struct examples. It exercises lower IR generation through
the smoke tests. It verifies that the build environment can compile the Rust
compiler.

\section{Performance}

The thesis does not present a detailed performance benchmark. The compiler is
not optimized for compile-time speed or generated-code speed. The most important
performance-related facts are architectural:

\begin{itemize}
  \item compile-time evaluation currently searches and replaces \texttt{comptime}
        expressions in a simple repeated traversal;
  \item lower IR represents local variables as stack allocation slots instead of
        promoting them to SSA values;
  \item LLVM object emission is performed without an enabled optimization pass
        pipeline;
  \item tree and CFG cleanup are intentionally local and conservative.
\end{itemize}

This is acceptable for the project scope. The compiler is designed to be
understandable and complete, not fast. Future work can add performance
benchmarks once language semantics and test coverage are stronger.

\section{Current Limitations}

The main limitations of \epl{} are:

\begin{itemize}
  \item one source file is compiled at a time;
  \item there are no modules, imports, packages, or separate compilation;
  \item arrays require literal lengths in type expressions;
  \item there are no generics;
  \item there is no standard library;
  \item pointer operations are unsafe and unchecked;
  \item array indexing has no bounds checks;
  \item the C ABI is only partially modeled;
  \item diagnostics stop at the first error and do not recover;
  \item lower IR optimizations are minimal;
  \item compile-time evaluation cannot use pointers or strings;
  \item there is no automatic final link command.
\end{itemize}

These limitations are expected for a final year compiler project. The important
point is that they are explicit. The implemented compiler has a coherent subset
rather than a large number of partially designed features.

\section{Future Work}

The first future improvement is stronger testing. The project should add
negative tests for parser errors, type errors, purity violations, invalid casts,
invalid field accesses, invalid array indexing, bad \texttt{main} signatures,
and malformed function declarations. The harness could also compare expected
diagnostic output.

The second improvement is run-time output testing. The compiler already emits
object files. A test runner could compile selected examples, link them in a
temporary directory, run the executables, and compare stdout and exit status.
This would test the LLVM backend and FFI path more directly than lower IR smoke
tests.

The third improvement is modules. Even a simple import system would require
decisions about file discovery, duplicate definitions, visibility, cyclic
dependencies, and separate or whole-program compilation. Modules would make the
language more usable but would significantly expand the compiler architecture.

The fourth improvement is a more complete ABI model. Reliable C interop requires
platform-specific rules for aggregate passing, varargs, return values, alignment,
and symbol naming. This would make \epl{} more practical as a systems language.

The fifth improvement is safety analysis. Bounds checks could be added for
arrays, and pointer operations could be made more explicit or restricted.
Implementing ownership or borrow checking would be a much larger research
direction, but smaller safety steps are possible.

The sixth improvement is optimization. Local alloca promotion, constant
propagation, dead-store elimination, common subexpression elimination, and LLVM
optimization pipelines could improve output quality. The current lower IR is a
good starting point because it already has basic blocks and immutable
definition IDs.

The seventh improvement is richer compile-time evaluation. It could support
constant expressions in array lengths, more built-in compile-time functions,
compile-time strings, or a separate model for compile-time-only references.
Each addition should preserve the purity boundary so that compile-time execution
remains deterministic.

\section{Conclusion}

This thesis presented the design and implementation of \epl{}, a small
ahead-of-time compiled language with compile-time evaluation. The compiler
implements a complete path from source text to native object code: lexing,
parsing, semantic analysis, typed \irtree{} construction, purity checking,
compile-time interpretation, lower CFG IR generation, cleanup passes, LLVM IR
generation, and object emission.

The project demonstrates that a compact compiler can still include modern
language ideas such as expression-oriented control flow and compile-time
execution. The implementation is intentionally conservative, but it is complete
enough to compile examples involving recursion, loops, arrays, structs,
pointers, external calls, and non-trivial compile-time computation. The current
test harness verifies the examples and regression tests with 24 passing checks.

The result is an educational compiler and language specification. It is not a
production language, but it is a solid foundation for further work in modules,
diagnostics, ABI support, safety, optimization, and richer compile-time
metaprogramming.
